
\begin{equation}\tag{5.1.1}
 i:\Pic(X)\longrightarrow K^{\bullet}(X)
\end{equation}
l'application qui associe à la classe \(\cl_P(\underline{L})\in\Pic(X)\) d'un \(\underline{O}_X\)-module inversible \(\underline{L}\) l'élément \(\cl^{\bullet}(\underline{L})\) de \(K^{\bullet}(X)\). Cette application est évidemment un homomorphisme de \(\Pic(X)\) dans le groupe multiplicatif \(K^{\bullet}(X)^*\) des éléments inversibles de \(K^{\bullet}(X)\). Par la suite nous noterons multiplicativement la loi de groupe de \(\Pic(X)\).

L'application \(i\) a un inverse à gauche défini par le foncteur «~det~» : Lorsque \(\underline{E}\) est un \(\underline{O}_X\)-module localement libre de rang \(\underline{n}\), on a
\[
 \det(\underline{E})=\bigwedge^{\underline{n}} \underline{E}.
\]
Si
\begin{equation}\tag{5.1.2}
0\longrightarrow \underline{E}'\longrightarrow \underline{E}\longrightarrow \underline{E}''\longrightarrow0
\end{equation}
est une suite exacte de \(\underline{O}_X\)-modules localement libres, il y a un isomorphisme canonique
\[
 \det(\underline{E})\xrightarrow{\sim}(\det(\underline{E}'))\otimes_{\underline{O}_X}(\det(\underline{E}''))
\]
provenant de (6.2.9) et du fait que (5.1.2) est localement scindée. On obtient donc un homomorphisme de groupes abéliens
\begin{equation}\tag{5.1.3}
 \det:K^{\bullet}(X)\longrightarrow\Pic(X)
\end{equation}
qui est l'inverse à gauche cherché de \(i\). En particulier \(i\) est injective, et nous identifierons par la suite \(\Pic(X)\) avec son image dans \(K^{\bullet}(X)^*\).

\subsection*{5.2.}
On définit une autre application
\begin{equation}\tag{5.2.1}
 \delta:\Pic(X)\longrightarrow\Filt^1(X)
\end{equation}
